\section{02 --- Datos, esquema canónico y preprocesado}

\subsection{1) Datasets y su papel}

\subsubsection{CICIDS2017 (principal)}

\begin{itemize}
  \item es la base moderna del proyecto
  \item se usan sus 8 CSV oficiales (\verb|src/load_cicids2017.py|)
  \item define el espacio de observación actual
\end{itemize}

\paragraph{Versiones del dataset}

Existen dos versiones locales:

\begin{table}[htbp]
\footnotesize
\begin{tabular}{llll}
\hline
\textbf{Versión} & \textbf{Ruta} & \textbf{Trackeada en git} & \textbf{Descripción} \\
\hline
Curada
  & \verb|datasets/CICIDS2017/*.csv|
  & Sí
  & Columnas con riesgo de leakage o redundantes eliminadas \\
  & & & antes de la ingesta. Es la que carga el adaptador. \\
Raw
  & \verb|datasets/CICIDS2017/Raw_dataset/|
  & No (en \texttt{.gitignore})
  & Exports CSV originales de CICFlowMeter. \\
  & & & Todas las columnas originales preservadas. Referencia local. \\
\hline
\end{tabular}
\end{table}

El adaptador (\verb|src/load_cicids2017.py|) aplica limpieza adicional en tiempo de carga
independientemente de qué versión se use. La política anti-leakage en código es la fuente de verdad
autoritativa. La versión curada reduce la superficie de ingesta pero no reemplaza las protecciones
en código.

\subsubsection{NSL-KDD (histórico)}

\begin{itemize}
  \item sirve como benchmark histórico
  \item mantiene valor didáctico para evolución del proyecto
  \item su mapeo al esquema canónico es muy parcial (\texttt{NSL\_KDD\_TO\_CANON})
\end{itemize}

\subsection{2) Decisión arquitectónica clave: esquema canónico}

Archivo: \verb|src/canonical_schema.py|

Se define una lista fija \texttt{FEATURES\_CANON} de \textbf{76 features flow-based}. Esto garantiza
que cada posición del vector tenga siempre el mismo significado.

\subsubsection{Observación final}

\begin{itemize}
  \item 76 valores de features
  \item 76 valores de máscara de missingness
  \item total: \textbf{152 dimensiones}
\end{itemize}

\subsection{3) Máscara de missingness (semántica)}

\begin{itemize}
  \item \texttt{1 = presente/válida}
  \item \texttt{0 = ausente o imputada}
\end{itemize}

Esto ayuda a no confundir ``dato real'' con ``dato rellenado''.

\subsection{4) Política anti-leakage}

\verb|load_cicids2017.py| elimina columnas propensas a leakage:

\begin{itemize}
  \item IPs
  \item timestamps
  \item Flow ID
  \item puertos directos como proxy de etiqueta
\end{itemize}

Razonamiento: si el modelo aprende atajos espurios, las métricas offline pueden inflarse
artificialmente y fallar en generalización real.

\subsection{5) Preprocesado en CICIDS2017}

Pipeline principal en \verb|load_cicids2017.py|:

\begin{enumerate}
  \item carga por chunks
  \item localización robusta de columna de etiqueta
  \item coerción numérica
  \item \(\inf \to \text{NaN} \to \texttt{fillna}(0)\) en features
  \item mapeo a canónico + máscara
  \item split (\texttt{random} o \texttt{day})
  \item \texttt{StandardScaler} (fit en train, transform en test)
\end{enumerate}

\subsection{6) Modos de split y por qué importan}

\subsubsection{\texttt{random}}

\begin{itemize}
  \item estratificado 80/20
  \item útil para iterar rápido
  \item puede sobreestimar generalización
\end{itemize}

\subsubsection{\texttt{day}}

\begin{itemize}
  \item train y test en grupos de días/CSV distintos
  \item más exigente y más cercano a escenario real
\end{itemize}

\subsubsection{\texttt{exact CSV}}

\begin{itemize}
  \item usado en leave-one-exact-CSV-out
  \item deja fuera un CSV real completo por fold
\end{itemize}

\subsection{7) Puntos finos que conviene dominar}

\begin{itemize}
  \item En CICIDS2017, los NaN se rellenan antes de mapear; la máscara captura sobre todo
        disponibilidad/validez tras mapping, no ``histórico crudo'' completo de missingness por
        celda.
  \item El \texttt{StandardScaler} se aplica a toda la observación (incluida máscara), así que la
        red recibe máscara escalada, no bits puros 0/1.
\end{itemize}
